\chapter{Security, Privacy and Reliability}\label{ch:security}

The system was reviewed twice for security: a first pass during development and a second pass before the MVP
submission (\code{docs/security/mvp-security-audit.md}). Both were practical reviews by the AI assistant used by the
team---code reading, the test suites, targeted probes against the running stack and a check of the host's open
ports. \textbf{They were not penetration tests or a comprehensive security audit.} The threat model is a single user
on one machine with all application ports bound to \code{127.0.0.1}.

\section{Implemented Controls}

\begin{table}[htbp]
  \centering
  \caption{Security and privacy controls and how they were verified}\label{tab:controls}
  {\small
  \begin{tabularx}{\textwidth}{@{}lLL@{}}
    \toprule
    \textbf{Area} & \textbf{Control} & \textbf{Verification} \\
    \midrule
    Upload validation & 50\,MB cap (nginx 55\,MB), PDF signature, full open, encryption and page checks before any job & Tests; live probes: non-PDF, truncated and 51\,MB files 422, 60\,MB 413 \\
    Storage & Server-generated \code{<sha256>.pdf} names; keys validated to stay inside the storage root & Test of unsafe keys; upload named \code{../../etc/passwd.pdf} stored inside storage \\
    Database & ORM and bound parameters only; lexical query terms reduced to word characters & Code search; injection-style queries left tables intact \\
    Outbound requests & arXiv hosts allow-listed (HTTPS), redirects re-checked, size caps, \code{defusedxml} & Tests \\
    Prompt injection & Delimited, neutralised passages; instructions to ignore passage instructions; JSON schema & Real-model test \\
    Citation integrity & Labels only; server-side resolution; answers without valid citations withheld & Tests; database-wide checks (\autoref{ch:testing}) \\
    Errors & Typed errors to status codes; fixed message for unexpected errors; 503 for database outage & Tests; live probes \\
    Resource bounds & Limits on query and question length, top-$k$, pagination, context, tokens, timeouts, one ingestion worker & Tests \\
    CORS and headers & Origin allow-list; CSP, \code{nosniff}, frame denial, no referrer & Test; \code{curl -I} \\
    Secrets & \code{.env} ignored by git; only \code{.env.example} committed; settings endpoint without credentials & Git history scan; test \\
    Privacy & No hosted model; logs contain event names, identifiers and counts, not questions or document text & Log review \\
    Dependencies & \code{pip-audit} and \code{npm audit} in CI & No known vulnerabilities at the time of review \\
    \bottomrule
  \end{tabularx}
  }
\end{table}

\section{Untrusted Retrieved Text and Model Output}

Retrieved passages are treated as untrusted data (\autoref{ch:litreview}). The model output is also untrusted: it
must validate against a schema, its citations are resolved against the supplied passages, and a failure produces an
explicit error rather than an answer. A structurally valid citation is labelled \emph{cited}, not \emph{supported},
and the interface states that semantic support is not verified.

\section{Local Model Deployment}

Generation runs only on the user's machine through Ollama; there is no hosted fallback. During the second review,
the host's Ollama service was found listening on all network interfaces, with inbound firewall rules allowing access
on a network classified as public. This exposes Ollama's unauthenticated API to the local network. It is a host
setting rather than an application defect; the remediation (binding Ollama to \code{127.0.0.1} and disabling the
inbound rules) is documented and must be applied by the machine owner. It was not applied during the review.

\section{Logging and Retention}

Logs are JSON lines with request identifiers. The web server's access log, which recorded query strings, was
disabled. The database retains papers, chunks and the full question history, including evidence snapshots, until
the user deletes them; there is no automatic retention limit.

\section{Availability and Recovery}

Readiness reports each dependency separately; search continues without Ollama. Jobs interrupted by a restart are
marked failed on start-up and can be retried. Availability was measured only over a 25.5-minute window (52 of 52
probes successful), which is too short to support the 99\% target. Backups use \code{pg\_dump} together with a copy
of the storage volume (\code{docs/operations.md}).

\section{Open Risks}

\begin{itemize}
  \item No authentication or authorisation: acceptable only while the system stays single-user and local.
  \item The application connects to PostgreSQL as a superuser; a least-privilege runtime role is planned.
  \item PDF parsing has no wall-clock timeout (bounded by the size and page limits and a single worker).
  \item Model revisions are pinned by name, not by revision, and are downloaded on first start.
  \item The CPU build of PyTorch is not covered by \code{pip-audit}.
\end{itemize}
